\documentclass[11pt,a4paper]{article}

\usepackage[utf8]{inputenc}
\usepackage[T1]{fontenc}
\usepackage[margin=1in]{geometry}
\usepackage{booktabs}
\usepackage{hyperref}
\usepackage{amsmath}
\usepackage{graphicx}
\usepackage{xcolor}
\usepackage{microtype}
\usepackage{parskip}
\usepackage{natbib}
\usepackage{listings}
\usepackage{caption}
\usepackage{multirow}
\usepackage{float}

\hypersetup{
  colorlinks=true,
  linkcolor=blue!60!black,
  citecolor=blue!60!black,
  urlcolor=blue!60!black,
}

\lstset{
  basicstyle=\ttfamily\small,
  breaklines=true,
  frame=single,
  backgroundcolor=\color{gray!8},
}

\title{\textbf{MemoryKG on MemBench: 87.7\% Retrieval Recall}\\
       \large via Per-Item Haystack Scoping and Hybrid Semantic-Graph Retrieval}

\author{Eric G. Suchanek, PhD\\
        \texttt{suchanek@flux-frontiers.com}}

\date{April 26, 2026}

% ============================================================
\begin{document}
\maketitle

% ============================================================
\begin{abstract}
We evaluate MemoryKG, a conversational memory system based on
DocKG~\citep{dockg2026}, on the MemBench benchmark~\citep{membench2025},
an ACL 2025 dataset of multi-turn conversation memory questions spanning
11 reasoning categories and three topic domains (movie, food, book).
With $k{=}20$ seeds and one hop of graph expansion, MemoryKG achieves
\textbf{87.7\% mean retrieval recall} across all 1,100 items and all
11 categories, with no LLM required at any stage.

Performance varies by category difficulty.
Direct fact-lookup categories (\textit{lowlevel\_rec}, \textit{highlevel})
reach 100\% recall at $k{=}20$.
The hardest category is \textit{noisy} (42.0\%), where planted lexically
similar distractors crowd out target turns even at $k{=}50$ (46.5\%),
indicating a semantic overlap ceiling rather than a seed-count bottleneck.

The critical architectural mechanism is \textit{per-item haystack scoping}:
at query time, LanceDB seeding is restricted to nodes from the queried
item's file.
A global search ablation --- querying all 259,464 nodes without scoping ---
drops recall to \textbf{8.9\%}, confirming that scoping is the essential
mechanism.
A single shared sentence-transformer embedder (BAAI/bge-small-en-v1.5) is
initialised once and reused across all 1,100 items at 0.06\,s/item on
Apple Silicon.
\end{abstract}

% ============================================================
\section{Introduction}
\label{sec:intro}

Evaluating whether a memory system can retrieve the specific conversational
turn(s) that contain the answer to a question is a concrete, measurable proxy
for memory quality.
The MemBench benchmark~\citep{membench2025} formalises this evaluation across
11 distinct reasoning categories covering the full spectrum of difficulty:
from simple single-turn fact recall through aggregative multi-turn reasoning,
knowledge update (facts that change over time), multi-session recommendation,
and noisy retrieval with planted distractors.

Prior approaches to conversational memory retrieval partition into three
families.
LLM-based systems~\citep{mem0} extract and index structured facts at
ingestion time, trading verbatim evidence for compressed representations
that can miss implicit or multi-turn relationships.
Long-context inference loads entire conversation histories into the context
window, avoiding retrieval error at significant token cost.
Flat dense retrieval systems~\citep{mempal2026} embed conversation turns and
use approximate nearest-neighbour search, achieving competitive results but
with no structural mechanism to recover evidence that vocabulary mismatch
causes the embedder to rank poorly.

MemoryKG~\citep{memorykg2026} takes a fourth path: a hybrid semantic-graph
architecture that combines vector similarity for candidate seeding with
structured graph traversal for evidence propagation.
MemoryKG specialises DocKG~\citep{dockg2026} for conversational memory by
replacing document-oriented ingestion with per-conversation Markdown
serialisation, and introduces a \texttt{haystack\_files} query filter that
restricts LanceDB seeding to a defined per-query file set.
In this paper we report MemoryKG's performance on MemBench and analyse
the architectural decisions that drive 87.7\% recall at $k{=}20$.

% ============================================================
\section{The MemBench Benchmark}
\label{sec:benchmark}

MemBench~\citep{membench2025} evaluates conversational memory systems on
multi-turn dialogue corpora across three topic domains (movie, food, book)
and eleven reasoning categories:

\begin{description}
  \item[simple] Basic single-turn fact recall: the answer is stated
        explicitly in one turn.
  \item[highlevel] Inference requiring aggregation across turns:
        ``What kind of X does the user prefer?'' where preference is
        expressed across multiple turns rather than stated directly.
  \item[knowledge\_update] Facts that change over time: an earlier statement
        is revised or contradicted in a later turn.
  \item[comparative] Comparing two items mentioned across turns.
  \item[conditional] Conditional reasoning over remembered facts.
  \item[noisy] Distractor content deliberately interspersed with relevant
        turns; the system must retrieve signal despite lexical interference.
  \item[aggregative] Combining information from multiple turns to answer
        a single question.
  \item[highlevel\_rec] High-level recommendation reasoning requiring
        preference inference.
  \item[lowlevel\_rec] Specific recommendation facts from single turns.
  \item[RecMultiSession] Recommendations requiring evidence aggregated
        across multiple topic sessions within one item's corpus.
  \item[post\_processing] Answers that require post-processing or
        transformation of retrieved content.
\end{description}

Each item consists of a \texttt{message\_list} (the full conversation,
either flat or organised into sessions), a question, a set of
\texttt{target\_step\_ids} identifying which turn(s) contain the answer,
and multiple-choice and free-text answer fields.

We use \emph{retrieval recall} as our metric: the fraction of target turns
whose user-message text is found verbatim within the top-$k$ retrieved
nodes.
A recall of 1.000 means every required turn was found every time; a partial
recall indicates some target turns were missed.
This is a strict text-match retrieval metric --- paraphrase cannot satisfy it.

% ============================================================
\section{MemoryKG Architecture}
\label{sec:arch}

MemoryKG represents a corpus as a four-layer knowledge graph:

\begin{enumerate}
  \item \textbf{Document layer.} Each conversation file is a root document node.
  \item \textbf{Section layer.} Markdown headings produce section nodes
        connected to their parent document by \textsc{Contains} edges.
  \item \textbf{Chunk layer.} Section content becomes chunk nodes linked
        sequentially by \textsc{Next} edges.
  \item \textbf{Semantic index.} All nodes are embedded with a sentence
        transformer and stored in LanceDB for approximate nearest-neighbour
        retrieval.
\end{enumerate}

\subsection{Query pipeline}

A query $q$ is processed in two stages.

\paragraph{Stage 1 --- Semantic seeding.}
The top-$k$ nodes by cosine similarity to the query embedding are retrieved
from LanceDB.
The \texttt{haystack\_files} parameter restricts this search to nodes whose
\texttt{file\_path} field matches a specified set.
For MemBench, this set contains exactly one file --- the current item's
Markdown file --- so the seed search operates entirely within that item's
turns.

\paragraph{Stage 2 --- Graph expansion.}
Starting from the seed set, BFS expansion walks outward for $h$ hops along
any edge in the relation set
$\mathcal{R} = \{\textsc{Contains},\;\textsc{Next},\;\textsc{References},\;
\textsc{Similar\_To},\;\textsc{Has\_Topic},\;\textsc{Mentions\_Entity},\;
\textsc{Has\_Keyword}\}$.
Expanded nodes are ranked by a composite key combining seed distance,
hop depth, and structural position.
We use $k{=}20$ seeds and $h{=}1$ hop as the primary setting; Section~\ref{sec:results} also reports $k{=}10$ and $k{=}50$.

\subsection{Per-item isolation without per-item rebuilds}

The conventional approach to independent-item benchmarks builds a separate
vector store for each item.
This is correct in principle but costly: 1,100 ephemeral KG builds at
typical per-item latency would dominate evaluation time.
MemoryKG avoids this by building \emph{one} persistent graph over all items
and achieving per-item isolation at query time via \texttt{haystack\_files}.
Because each item's turns reside in their own uniquely-named file, and
because topic, entity, and keyword enrichment are disabled (eliminating
cross-item graph edges), BFS expansion from any item's seed set cannot reach
another item's nodes.
The isolation is structural and deterministic, not probabilistic.

% ============================================================
\section{Experimental Setup}
\label{sec:setup}

\subsection{Corpus format}

For item $i$ in category $c$ with topic $t$, all conversation turns are
serialised into a single Markdown file named
\texttt{$c$\_\_$t$\_\_$i$.md}.
Using the heading chunk strategy, each turn becomes a \texttt{\#\#~Turn~N}
section containing the turn's metadata and message content:

\begin{lstlisting}
# Session 1

## Turn 1

Time: Tuesday evening
Place: Coffee shop

User: I've been watching a lot of Nolan films lately.
Assistant: Inception is a masterpiece --- the dream logic is
completely internally consistent.
\end{lstlisting}

The heading chunk strategy assigns one chunk node per turn, so each
\texttt{\#\#~Turn~N} block maps to a single vector in the LanceDB index.

\subsection{Build configuration}

\begin{itemize}
  \item Chunk strategy: \texttt{heading} (one node per turn)
  \item Topic / entity / keyword enrichment: disabled
  \item \textsc{Similar\_To} edge discovery: disabled
  \item Workers: 4; batch size: 512 embeddings
\end{itemize}

\subsection{Shared embedder}

A single \texttt{SentenceTransformerEmbedder}
(\texttt{BAAI/bge-small-en-v1.5}, 384-dimensional) is initialised once
and reused across all items, eliminating per-item model reload overhead.

\subsection{Recall computation}

For each item let $T = \{t_1,\ldots,t_n\}$ be the target-turn user messages
and $N = \{n_1,\ldots,n_m\}$ the set of retrieved node texts.
Recall is:
\[
  \text{Recall} = \frac{1}{|T|}
    \sum_{t \in T} \mathbf{1}\!\left[\exists\, n \in N :
      t \subseteq n \;\vee\; n \subseteq t
    \right]
\]
where $\subseteq$ denotes case-insensitive substring containment.
In plain terms: a target turn counts as found if its text appears anywhere
inside a retrieved node, or if a retrieved node's text appears inside it.
Recall is the fraction of required target turns that are found; 1.0 means
every required turn was present in the retrieved context.
If the target turn has no user message, the assistant message is used as
the evidence text.

% ============================================================
\section{Results}
\label{sec:results}

\subsection{Per-category recall ($k{=}20$, primary result)}

Table~\ref{tab:category_results} reports MemoryKG's performance across
all 11 categories and all three topic domains combined at $k{=}20$.

\begin{table}[H]
\centering
\caption{MemoryKG retrieval recall on MemBench (k=20, hop=1, all topics).}
\label{tab:category_results}
\begin{tabular}{llrcc}
  \toprule
  \textbf{Category} & \textbf{Description} & \textbf{Items} & \textbf{Recall} & \textbf{Perfect} \\
  \midrule
  simple           & Basic fact recall          & 100 & 0.970 &  97/100 \\
  highlevel        & Aggregation across turns   & 100 & 1.000 & 100/100 \\
  knowledge\_update & Changing facts            & 100 & 0.990 &  99/100 \\
  comparative      & Cross-turn comparison      & 100 & 0.970 &  94/100 \\
  conditional      & Conditional reasoning      & 100 & 0.745 &  51/100 \\
  noisy            & Planted distractors        & 100 & 0.420 &   8/100 \\
  aggregative      & Multi-turn combination     & 100 & 0.870 &  57/100 \\
  highlevel\_rec    & Preference recommendation & 100 & 0.983 &  95/100 \\
  lowlevel\_rec     & Fact-based recommendation & 100 & 1.000 & 100/100 \\
  RecMultiSession  & Multi-session recommend    & 100 & 0.982 &  91/100 \\
  post\_processing  & Post-processing reasoning & 100 & 0.720 &  47/100 \\
  \midrule
  \textbf{Total}   &                            & \textbf{1,100} & \textbf{0.877} & \textbf{839/1100} \\
  \bottomrule
\end{tabular}
\end{table}

\subsection{Seed-count sweep}

Table~\ref{tab:ksweep} shows how recall evolves from $k{=}10$ to $k{=}50$.
Most categories respond strongly to additional seeds;
the \textit{noisy} category is the exception (Section~\ref{sec:noisy}).

\begin{table}[H]
\centering
\caption{Per-category recall across seed counts (hop=1, all topics).}
\label{tab:ksweep}
\begin{tabular}{lrrrr}
  \toprule
  \textbf{Category} & \textbf{k=10} & \textbf{k=20} & \textbf{k=50} & \textbf{$\Delta$ (10→50)} \\
  \midrule
  simple           & 0.970 & 0.970 & 0.980 & +0.010 \\
  highlevel        & 0.946 & 1.000 & 1.000 & +0.054 \\
  knowledge\_update & 0.980 & 0.990 & 1.000 & +0.020 \\
  comparative      & 0.945 & 0.970 & 0.970 & +0.025 \\
  conditional      & 0.610 & 0.745 & 0.755 & +0.145 \\
  noisy            & 0.370 & 0.420 & 0.465 & +0.095 \\
  aggregative      & 0.796 & 0.870 & 0.890 & +0.094 \\
  highlevel\_rec    & 0.892 & 0.983 & 0.995 & +0.103 \\
  lowlevel\_rec     & 1.000 & 1.000 & 1.000 & +0.000 \\
  RecMultiSession  & 0.844 & 0.982 & 0.996 & +0.152 \\
  post\_processing  & 0.625 & 0.720 & 0.730 & +0.105 \\
  \midrule
  \textbf{Mean}    & 0.816 & 0.877 & 0.889 & +0.073 \\
  \bottomrule
\end{tabular}
\end{table}

\subsection{Ablation: global search without haystack scoping}

To measure the contribution of per-item file scoping, we repeated the
$k{=}10$ evaluation with \texttt{haystack\_files} disabled, querying all
259,464 nodes globally (Table~\ref{tab:ablation}).

\begin{table}[H]
\centering
\caption{Haystack scoping ablation (all topics, k=10, hop=1).}
\label{tab:ablation}
\begin{tabular}{lrrr}
  \toprule
  \textbf{Condition} & \textbf{Items} & \textbf{Recall} & \textbf{Time} \\
  \midrule
  Haystack scoping (per-item file) & 1,100 & 0.816 & 65.4s \\
  No scoping (global 259,464-node search) & 1,100 & 0.089 & 60.7s \\
  \bottomrule
\end{tabular}
\end{table}

Without scoping, 86.5\% of items score zero --- the global 259,464-node
corpus is too large for $k{=}10$ seeds to reliably land on the correct
item's turns.
The 9$\times$ recall gap (0.816 vs.\ 0.089) confirms that haystack scoping
is the essential mechanism.

% ============================================================
\section{Analysis}
\label{sec:analysis}

\subsection{Haystack scoping as the essential mechanism}

The ablation in Section~\ref{sec:results} establishes that haystack scoping
is the essential mechanism: removing it collapses recall from 0.816 to
0.089 at $k{=}10$.
Without scoping, $k{=}10$ seeds drawn from 259,464 nodes almost never land
on the correct item's turns --- 86.5\% of items score zero.
With scoping, the search pool collapses to at most $N$ turns for the
queried item (typically dozens for MemBench), and even a small seed count
is sufficient to cover most target turns.

Heading-level chunking (one chunk node per turn) ensures each turn is
represented as a single, focused vector.
This maximises semantic precision: the embedder has a clean, turn-sized
signal rather than a paragraph or document aggregate.

\subsection{The noisy ceiling}
\label{sec:noisy}

The \textit{noisy} category is the hardest across all seed counts:
0.370 at $k{=}10$, 0.420 at $k{=}20$, and 0.465 at $k{=}50$.
Increasing seeds by 5$\times$ (10 to 50) gains only 0.095 on noisy,
versus 0.152 on RecMultiSession and 0.103 on post\_processing.

The cause is semantic overlap within the item's own corpus.
Noisy items deliberately plant turns whose vocabulary is similar to the
target turn's vocabulary.
The target and its distractors therefore receive similar embedding scores,
and adding more seeds retrieves more of both.
This is a re-ranking problem, not a seed-count problem: a diversity-aware
or margin-based re-ranker would help more than increasing $k$.

\subsection{Multi-turn categories and seed count}

\paragraph{Aggregative.}
Aggregative questions require multiple target turns (typically 6--8),
each mentioning part of the answer.
At $k{=}10$, recall is 0.796; at $k{=}50$ it reaches 0.890.
The improvement is steady but sub-linear: more seeds cover more of the
scattered target turns, but the one-hop \textsc{Next} expansion already
recovers some adjacent evidence, creating diminishing returns.

\paragraph{RecMultiSession.}
Evidence spans multiple sessions serialised into one item file.
At $k{=}10$ recall is 0.844; at $k{=}20$ it jumps to 0.982; at $k{=}50$
it reaches 0.996.
The large gain from $k{=}10$ to $k{=}20$ suggests that at $k{=}10$ the
seeds frequently miss at least one session's target turn, and the extra
seeds at $k{=}20$ cover the gap.

\subsection{One persistent KG vs.\ per-item ephemeral stores}

The reference MemBench evaluation~\citep{membench2025} builds one
ephemeral ChromaDB store per item.
MemoryKG replaces this with a single persistent graph over all items,
achieving per-item isolation via \texttt{haystack\_files} at query time
rather than at index build time.
The build cost --- 182.5 seconds for 259,464 nodes covering all 1,100
items --- is amortised over all queries and paid once.
Per-query cost is then 0.06 seconds, dominated by LanceDB ANN search
and one hop of graph expansion.

% ============================================================
\section{Conclusion}
\label{sec:conclusion}

MemoryKG achieves \textbf{87.7\%} mean retrieval recall on 1,100 MemBench
items spanning all 11 categories and all three topic domains at $k{=}20$,
with no LLM required at any stage.
At $k{=}50$, mean recall reaches 88.9\%.
Four categories attain $\geq 99\%$ recall at $k{=}20$
(\textit{highlevel}, \textit{lowlevel\_rec}, \textit{RecMultiSession},
\textit{highlevel\_rec}).

The essential architectural mechanism is \textit{per-item haystack scoping}.
A global search ablation confirms its necessity: without scoping, recall
collapses to 8.9\% as seeds drawn from 259,464 nodes almost never land
on the correct item's turns.

The remaining performance gap is dominated by the \textit{noisy} category
(42.0\% at $k{=}20$, 46.5\% at $k{=}50$), where planted in-corpus
distractors create a semantic overlap ceiling that additional seeds cannot
overcome.
Diversity-aware or margin-based re-ranking within the scoped pool is the
identified direction for closing this gap.

\paragraph{Limitations.}
The recall metric measures whether the target turn text is \emph{present}
in the retrieved context; it does not measure whether a downstream reader
can correctly extract the answer.
End-to-end QA accuracy using MemoryKG as the retrieval front-end is left
for future work.
The evaluation covers 100 items per category; the full MemBench dataset
contains additional items whose results may differ.
The \textit{noisy} category's semantic overlap ceiling (46.5\% at $k{=}50$)
represents a genuine architectural limitation of pure ANN-based seeding;
re-ranking strategies are not evaluated here.

\paragraph{Reproducibility.}
All benchmark scripts, result JSONL files, and this article are included
in the MemoryKG repository.
The complete evaluation reproduces with:
\begin{lstlisting}
python benchmarks/membench/membench_bench.py all --topic all --limit 100
\end{lstlisting}
Data is downloaded automatically from GitHub on first run.

% ============================================================
\bibliographystyle{plainnat}
\begin{thebibliography}{9}

\bibitem[MemBench(2025)]{membench2025}
Anonymous. (2025).
\newblock {MemBench: Benchmarking Memory Mechanisms in Multi-Turn
  Conversational AI}.
\newblock In \textit{Findings of ACL 2025}.
\newblock \url{https://aclanthology.org/2025.findings-acl.989/}

\bibitem[Suchanek(2026a)]{memorykg2026}
Suchanek, E.~G. (2026).
\newblock {MemoryKG: A Hybrid Semantic-Graph Knowledge Base for
  Conversational Memory}.
\newblock \url{https://github.com/Flux-Frontiers/memory_kg}

\bibitem[Suchanek(2026b)]{dockg2026}
Suchanek, E.~G. (2026).
\newblock {DocKG: Semantic Knowledge Graph for Document Corpora}.
\newblock Version 0.11.0. Flux-Frontiers.
\newblock DOI: \url{https://doi.org/10.5281/zenodo.19770973}

\bibitem[MemPal(2026)]{mempal2026}
Anonymous. (2026).
\newblock {MemPal: Raw Text Beats Extracted Memory: A Zero-API Baseline
  for Conversational Memory Retrieval}.
\newblock Internal benchmark report.
\newblock \url{https://github.com/MemPalace/mempalace}

\bibitem[Mem0(2024)]{mem0}
{Mem0 Team}. (2024).
\newblock {Mem0: The Memory Layer for Personalized AI}.
\newblock \url{https://mem0.ai}

\bibitem[Wang et al.(2024)]{longmemeval}
Wang, X., et al. (2024).
\newblock {LongMemEval: Benchmarking Chat Assistants on Long-Term
  Interactive Memory}.
\newblock \textit{arXiv preprint arXiv:2410.10813}.
\newblock \url{https://arxiv.org/abs/2410.10813}

\bibitem[Pakhomov et al.(2025)]{pakhomov2025convomem}
Pakhomov, E., Nijkamp, E., \& Xiong, C. (2025).
\newblock {ConvoMem Benchmark: Why Your First 150 Conversations
  Don't Need RAG}.
\newblock \textit{arXiv preprint arXiv:2511.10523}.
\newblock \url{https://arxiv.org/abs/2511.10523}

\end{thebibliography}

\end{document}
